\chapter{Literature Review}
\label{Chapter 2}

\section{Previous Approaches to Solve the Problem}

\subsection{Existing Collection-Management Products}
If you look at what is actually in market in India today, the commercial collection-management space is owned by a handful of names whose products were built before the current wave of conversational AI got going. CreditMantri, Spocto, Cred -- those three plus the internal tools at a couple of the larger private banks cover most of what you'll see in real deployments. The shapes are all roughly the same: a case database underneath, some dialer integration on the side, reporting bolted on top. The thing genuinely missing in every one of them is a conversational layer that isn't just humans on phones. Where WhatsApp shows up, it shows up via Meta's Cloud API, which brings the usual baggage -- templates that need Meta's approval, throughput caps, per-message pricing that scales linearly. The academic side has very little written about collections software itself; the real work tends to happen inside banks and never gets out into a paper. What is published sits in the adjacent space. Verma, Mathur and Singh~\cite{vermaml} go through ML techniques applied to behavioural credit data, and a few of those carry over fairly naturally to picking the right channel per customer. Wang and Zhang~\cite{contactcenter} cover predictive dialers and the queue-management bits sitting under them. None of this prior work, though, looks at the specific combination we are aiming at here: AI voice plus AA-driven prioritisation plus consumer-SIM WhatsApp.

\subsection{Conversational Voice Agents}
The biggest single technology shift behind this kind of platform is the maturity of conversational voice agents. A modern voice agent is a pipeline -- automatic speech recognition (ASR), a large language model (LLM) for the actual dialogue, and a text-to-speech (TTS) synthesiser -- and each piece is now fast enough to meet real-time latency budgets that did not exist three years ago. The voice stack Ignosis uses in production is the firm's own internal voice AI platform~\cite{bolna}: it chains a low-latency ASR, a fine-tuned LLM and a TTS engine that produces Indian English with regional-accent variation. Its target latency is around 700~ms from the customer finishing speaking to the agent beginning to reply, which sits comfortably below the 800--900~ms band at which human listeners start to feel something is off. Open-source baselines such as Pipecat~\cite{pipecat} and LiveKit Agents~\cite{livekit} were evaluated briefly before the internal platform was settled on. The decision came down to two things: how well the ASR handled rural and semi-urban customers in Indian English, and how mature the post-call extraction was. Table~\ref{tab:voice-agents} captures the comparison.

\begin{table}[H]
\centering
\caption{Comparison of conversational voice-agent stacks evaluated for the platform.}
\label{tab:voice-agents}
\small
\setlength{\tabcolsep}{4pt}
\begin{tabularx}{\textwidth}{@{}>{\raggedright\arraybackslash}p{0.30\textwidth} >{\centering\arraybackslash}X >{\centering\arraybackslash}X >{\centering\arraybackslash}X@{}}
\toprule
\textbf{Capability} & \textbf{Internal voice AI platform} & \textbf{Pipecat} & \textbf{LiveKit Agents} \\
\midrule
Indian-English ASR accuracy            & High             & Medium         & Medium \\
Indic regional language support        & Yes (5+)         & Limited        & Limited \\
Median turn latency                    & $\sim$700~ms     & $\sim$900~ms   & $\sim$850~ms \\
Telephony (PSTN) integration           & Built-in         & DIY            & DIY \\
Post-call structured extraction        & Built-in         & Manual         & Manual \\
Webhook-based outcome delivery         & Built-in         & Custom         & Custom \\
Hosting model                          & Managed in-house & Self-hosted    & Self-hosted \\
Operational maturity for collections   & Production-ready & Framework-only & Framework-only \\
\bottomrule
\end{tabularx}
\end{table}

\subsection{The WhatsApp Multi-Device Protocol}
WhatsApp's multi-device protocol lets up to four secondary devices share end-to-end-encrypted access to a primary phone's account without that phone being online. The official integration path for businesses is the WhatsApp Cloud API, which exposes a REST interface and forces every first message in a conversation to be a pre-approved template. For collections, this is operationally restrictive in three ways. Every template must be approved by Meta, which can take days. The per-message price makes the channel uneconomic for low-ticket cases. And in practice customers respond much better to messages from real consumer-grade WhatsApp numbers than from obviously branded business accounts. To work around these, the open-source community has produced several reimplementations of the WhatsApp protocol; the project uses one of the most widely deployed of these~\cite{baileys}.

\subsection{Anti-Spam and Rate-Limiting Heuristics}
A real chunk of the engineering work on SIM-based WhatsApp goes into staying ahead of Meta's spam-detection heuristics. The exact rules are not published, but the open-source community has converged on a handful of practices that visibly reduce flagging. Mihail et al.~\cite{antispam} review the academic literature on spam detection in messaging platforms, and the same themes show up over and over: randomised intervals between messages, language-level variation to defeat hash-based fingerprinting, simulated typing presence, and conservative daily caps. The platform built in this project hits each of these.

\subsection{Account Aggregator Framework}
The AA framework~\cite{aarbi} is the RBI-regulated path through which a registered Financial Information User (FIU) can pull a customer's bank statements, with consent that is explicit, time-bound and revocable at any moment. Our collections platform pulls that data because it changes the AI script meaningfully -- the script can now branch on actual signals like a fresh salary credit, or a balance trending down, and what used to be a flat ``please pay'' opening becomes the much sharper ``we can see your salary credit landed on the 1st, what date works''.

\section{Problem Identification}
Reading the prior art against what it doesn't quite cover, the engineering work in this project breaks down into nine fairly concrete problems.

First, cycle-level case management -- ingesting the monthly delinquent CSV the lender hands over, validating it, de-duping it, allocating cases to agents per the tenant's own rules, and getting them onto a paginated dashboard. Second, the case detail screen itself, which has to gather everything about a customer onto a single page (contact info, loan and payment history, past call transcripts, WhatsApp and SMS threads, AA signals) without making the page take five seconds to load -- which means lazy loading and progressive rendering. Third, bulk campaign execution: a way to wrap a batch of cases as one schedulable thing built from a CSV, with pause and resume, and wired into both the voice engine and the messaging layer. Fourth, the speed dialer -- the moment an agent submits a disposition, the next case must already be sitting on the screen, no in-between state. Fifth, AI voice calling itself, triggerable both from inside a campaign and manually as a fallback, with the callback parsed and the workflow walked forward correctly.

Then the messaging side. Sixth, WhatsApp automation: a pool of consumer SIMs, rotation across them, per-message body variation, daily cap enforcement per sender, and automatic re-routing around banned numbers without anyone manually stepping in. Seventh, SMS automation: dispatch from a real consumer SIM on a real Android phone, driven by a small app. Eighth, auditability -- log everything densely enough that QA can later reconstruct what happened to any one case at any point in time. And ninth, multi-tenant isolation: more than one lender tenant on the same deployment, with strict separation of data, identity and configuration so nothing crosses over.

Those nine items are the working set Chapter~3 will be judged against.
